\documentclass[12pt,oneside,a4paper]{book}
\usepackage[latin1]{inputenc}
\usepackage[english, spanish]{babel}	%For languages characters and hyphenation
\usepackage{amsmath}
\usepackage{amsfonts}
\usepackage{titlesec}
\usepackage{amssymb}
\usepackage{graphicx}
\usepackage{array}
\usepackage{cite}
\usepackage[all]{xy}
\usepackage{fancyhdr}
\usepackage{lastpage}
\usepackage{caption}
\usepackage{subcaption}
\usepackage{microtype}		%Makes pdf look better.
\usepackage[left=3.00cm, right=3.00cm, top=2.00cm, bottom=2.00cm]{geometry}
\author{Julian David Mora Ramos}
\title{Proyecto, implementaci�n de SIGEPI}
\usepackage[pdftex,
	pdfauthor={JULIAN DAVID MR},
	pdftitle={IMPLEMENTACI�N DEL SISTEMA SIGEPI},
	pdfsubject={DESARROLLO DE SOFTWARE},
	pdfkeywords={SOFTWARE, MDD, MDE, UDLA, DOCUMENTACI�N, DESARROLLO},
	pdfproducer={Universidad de la Amazonia},
	pdfcreator={@anlijudavid}]{hyperref}

% -------------------------
\setcounter{tocdepth}{4}
\setcounter{secnumdepth}{4}
\newcommand{\myparagraph}[1]{\paragraph{#1}\mbox{}\\}	
% -------------------------

% ----------------Tama�o de texto Chapters-----------------
\newcommand{\chapfnt}{\fontsize{16}{19}}
\newcommand{\secfnt}{\fontsize{14}{17}}
\newcommand{\ssecfnt}{\fontsize{12}{14}}

\titleformat{\chapter}[display]
{\normalfont\chapfnt\bfseries}{\chaptertitlename\ \thechapter}{20pt}{\chapfnt}

\titleformat{\section}
{\normalfont\secfnt\bfseries}{\thesection}{1em}{}

\titleformat{\subsection}
{\normalfont\ssecfnt\bfseries}{\thesubsection}{1em}{}

\titlespacing*{\chapter} {0pt}{50pt}{40pt}
\titlespacing*{\section} {0pt}{3.5ex plus 1ex minus .2ex}{2.3ex plus .2ex}
\titlespacing*{\subsection} {0pt}{3.25ex plus 1ex minus .2ex}{1.5ex plus .2ex}
% --------------------------------------------------------

% -----------------Config de grafica MVC------------------
\usepackage{tikz}
\usetikzlibrary{arrows,positioning} 
\tikzset{
	%Define standard arrow tip
	>=stealth',
	%Define style for boxes
	punkt/.style={
		rectangle,
		rounded corners,
		draw=black, very thick,
		text width=6.5em,
		minimum height=2em,
		text centered},
	% Define arrow style
	pil/.style={
		->,
		thick,
		shorten <=2pt,
		shorten >=2pt,}
}
% --------------------------------------------------------

% ------------Permite renderizar muchas tablas------------
\maxdeadcycles=1000
% --------------------------------------------------------

\begin{document}
	\markboth{"SIGEPI"}{''Universidad de la Amazonia''}
	
	% Paginas iniciales
	\title{%
	\begin{figure}[htbp]
		\centering
		\includegraphics[width=50px]{resources/logoudla.jpg}
	\end{figure} \\ %
	Implementaci�n del sistema de informaci�n SIGEPI basado en herramientas de desarrollo dirigido por modelos \\
	\large Implementaci�n de la plataforma web para la gesti�n de los proyectos presentados por los grupos y semilleros de investigaci�n en la Universidad de la Amazonia \\
}
\author{
	Julian David Mora Ramos $^{1}$
	\\
	\small{$^{1}$Universidad de la Amazonia, $^{1}$Florencia,Caquet�}\\
}
\maketitle
	\thispagestyle{empty}

\begin{center}
	\includegraphics{./resources/logoudla.jpg}
\end{center}
\vspace{2cm}

\begin{center}
	{\Large {\bf UNIVERSIDAD DE LA AMAZONIA}}
	\vspace{5mm}

 	{\Large {Ingenier�a de Sistemas}}
    \\
  	\vspace{34mm}
	{\large {\bf PROYECTO FIN DE PREGRADO}}
  	\vspace{10mm}
    \\
  	{\Large {{\Huge {				
  		IMPLEMENTACI�N DEL SISTEMA INFORMACI�N PARA GESTI�N DOCUMENTAL SIGEPI
	}} \\[1cm] }}
  	\vspace{2cm}
	{\large {
  		{\bf Autor}:\\
  		Julian David Mora Ramos\\
  	}}
	\vspace{10mm}
  	{\large {Curso acad�mico 2011/2017}}
  	\vspace{1cm}
\end{center}
	%------------------------------------------------
%   DEDICATORIA
%------------------------------------------------

\pagestyle{empty}
\frontmatter

\chapter*{}
\begin{flushright}
	\textit{DEDICATORIA}
		
	A la Facultad de Ingenier�a y a la  Universidad, por la formaci�n que me han dado.\\
	Es gracias a ustedes que es posible el presente trabajo.\\
	En verdad, gracias.\\
	Yo.
\end{flushright}
	%------------------------------------------------------
%   AGRADECIMIENTOS
%------------------------------------------------------

\chapter*{Agradecimientos}
%\markboth{AGRADECIMIENTOS23}{AGRADECIMIENTOS} % encabezado 

�Muchas gracias a todos!

	\section*{Nota de aceptaci�n}
	
Tambi�n quisiera reconocer a ... por ...CONACYT,  PAPIIT / etc.



	%--------------------------------------------------
%   DECLARACI�N
%--------------------------------------------------

\thispagestyle{empty}
\vspace*{\fill}
\begingroup
``Con fundamento en los art�culos 21 y 27 de la Ley Federal del Derecho de Autor y como titular de los derechos moral y patrimonial de la obra titulada ``\textbf{T�TULO DE LA TESIS}'', otorgo de manera gratuita y permanente al Instituto Tecnol�gico Aut�nomo de M�xico y a la Biblioteca Ra�l Baill�res Jr., la autorizaci�n para que fijen la obra en cualquier medio, incluido el electr�nico, y la divulguen entre sus usuarios, profesores, estudiantes o terceras personas, sin que pueda percibir por tal divulgaci�n una contraprestaci�n''.

\centering

\hspace{3em}

\textsc{AUTOR}

\vspace{5em}

\rule[1em]{20em}{0.5pt} % L�nea para la fecha

\textsc{Fecha}
 
\vspace{8em}

\rule[1em]{20em}{0.5pt} % L�nea para la firma

\textsc{Firma}

\endgroup
\vspace*{\fill}

	\begin{center}	
	\section*{Resumen}
	
	
	\section*{Abstract}
\end{center}
	
	\tableofcontents
	
	\cleardoublepage
	\addcontentsline{toc}{chapter}{Lista de figuras} % para que aparezca en el indice de contenidos
	\listoffigures % indice de figuras
	
	\cleardoublepage
	\addcontentsline{toc}{chapter}{Lista de tablas} % para que aparezca en el indice de contenidos
	\listoftables % indice de tablas
	
	% \pagenumbering{arabic} % para empezar la numeraci�n con n�meros
	% Empieza la numeraci�n de las p�ginas	
	\mainmatter
		
	% Contenido
	\section[INTRODUCCI�N]{INTRODUCCI�N}

Actualmente, adaptarse a las necesidades del cliente es un
problema que aumenta con el transcurso del tiempo, las condiciones
de negocio cambian y consigo tambi�n la construcci�n de sistemas.
Consecutivamente, el mundo moderno exige gran demanda tecnol�gica
y adaptarse a esta linea de constantes transformaciones resulta
en una tarea compleja. Por esto, gracias a las investigaciones
realizadas en este campo, se han logrado avances en ciertas �reas
fundamentales que influyen en la puesta en marcha de proyectos.
\\

Muchos proyectos de desarrollo de software poseen gran demanda, y
m�s cuando est�n soportados por herramientas de automatizaci�n
de tareas. Los tiempos de ejecuci�n y elaboraci�n de proyectos se
reducen cuando herramientas DSL (Lenguaje de dominio especifico)
son implementadas, ofreciendo caracter�sticas de generaci�n de
c�digo reutilizable y componentes.
\\

El lenguaje de dominio especifico no es una tecnolog�a emergente,
las primeras apariciones fueron a mediados de los 80s. 
Actualmente se pueden encontrar ramificaciones DSL 
en diferentes aplicaciones, cada una con
funcionalidades especificas, tales como MDD (Desarrollo Dirigido
por Modelos), MDA (Arquitectura Dirigida por Modelos) y
MDE (Ingenieria Orientada a Modelos), siendo estos un conjunto de
recursos interoperables que permiten ser usados en el
an�lisis, modelado y construcci�n de proyectos
software. Por otro lado, han aumentado la cantidad de mecanismos
para agilizar la planeaci�n y ejecuci�n de proyectos, generalmente
todos estos siguen la misma filosof�a de las DSL, por ejemplo, se
pueden encontrar las ORM (Mapeo de Objetos Relacionales) para el
control de m�ltiples gestores
de bases de datos a partir de un solo lenguaje base. Del mismo 
modo, los ingenieros, programadores de computadoras, matem�ticos,
estad�sticos y dem�s, hacen uso de herramientas de lenguaje de
dominio 
especifico a partir de programas de computador como R (Lenguaje
estad�stico), interpretes 
de expresiones regulares, SQL (lenguaje de consulta estructurada), 
LINQ (Lenguaje de consulta integrada), entre otros. Como se puede
observar, un dominio espec�fico tiene un vocabulario especializado
para describir las cosas que son particulares a ese dominio.
\\

El sistema de gesti�n documental, SIGEPI, es una plataforma web 
robusta, creada gracias a la implementaci�n de herramientas 
de modelado y generaci�n de c�digo, DSL. En la primera secci�n de 
este documento se nombran todos aquellos procesos que se siguieron 
para la obtenci�n del c�digo fuente a partir de los esquemas 
dados al aplicar MDD (Desarrollo Dirigido por Modelos).
Posteriormente, se explica en detalle la arquitectura de desarrollo que fue
aplicada, incluyendo cada una de las librer�as y marcos de trabajo que 
fueron usados.
	\section[DESCRIPCI�N GENERAL DEL PROYECTO]{DESCRIPCI�N GENERAL DEL PROYECTO}

\subsection[Planteamiento del problema]{Planteamiento del problema}

\subsubsection[Contexto]{Contexto}
\subsubsection[Formulaci�n del problema]{Formulaci�n del problema}

\subsection[Justificaci�n]{Justificaci�n}

\subsection[Objetivos]{Objetivos}
\subsubsection[Objetivo general]{Objetivo general}
\subsubsection[Objetivo especifico]{Objetivo especifico}

	\section[MARCO REFERENCIAL]{MARCO REFERENCIAL}

\subsection[Marco te�rico]{Marco te�rico}

\subsubsection[Desarrollo Dirigido por Modelos. Conceptos]{Desarrollo Dirigido por Modelos. Conceptos}
El desarrollo dirigido por modelos (MDD) es un paradigma que resuelve
muchos inconvenientes en el desarrollo de software, la causa 
surgi� desde los inicios de la d�cada de los 60s, cuando se
introdujo el concepto de ''crisis del software'', originado por la complejidad y el costo requerido por las necesidades del cliente \cite{pons2010desarrollo}.
MDD enriquece el proceso de desarrollo a partir de herramientas 
especializadas, destaca una relaci�n entre lo abstracto y el c�digo
fuente. La abstracci�n est� representada por modelos, por tanto 
estos requieren de una o m�s fases de transformaci�n para 
finalmente obtener el c�digo [ver figura \ref{fig_transcode}].

\begin{displaymath}
	\xymatrix{
		*+<1cm>[F-,]{Modelo}\ar[r] 
		& *+<1cm>[o][F]{M2T/M2M}\ar[r]
		& ... \ar[r]
		& *+<1cm>[F-,]{Modelo}\ar[r]
		& *+<.5cm>[F]{Codigo}
	}
\end{displaymath}
\captionof{figure}{Transformaci�n de c�digo}
\label{fig_transcode}


\subsubsection[DSL, Lenguaje de dominio especifico]{Lenguaje de dominio especifico}
El lenguaje de dominio especifico eleva el nivel de abstracci�n m�s all�
que los lenguajes de programaci�n para especificar una soluci�n al
problema usando conceptos de dominio \cite{kelly2008domain}. 
Generalmente se le encuentra como una notaci�n gr�fica, cuyos
modelos resultan en un conjunto de elementos y relaciones entre si.

\begin{displaymath}
	\xymatrix{
		\underline{Requerimientos}\ar@/_1mm/[r] 
		& *+<1cm>[o][F]{\texttt{Abstracci�n}}\ar[r]
		& *+<1cm>[F-,]{\textbf{DSL}}
	}
\end{displaymath}
\captionof{figure}{Abstracci�n de DSL}
\label{fig_abs_dsl}

\subsubsection[Meta modelos]{Meta modelos}


\subsubsection[Arquitectura MVC]{Arquitectura MVC}

\subsubsection[Framework de desarrollo Ext.NET]{Framework de desarrollo Ext.NET}
	\chapter{Procesos para presentaci\'on de proyectos}

Vicerector\'ia de investigaciones de la Universidad de la
Amazonia cuenta 

	\chapter{Base de datos}

Diccionario de datos 

\begin{table}[ht]
	\caption{Tabla Usuario}
	\label{labelTableUsuario}
	\begin{tabular}{ |l|l|l|l|l|l| }
		\hline
		Nombre & Tipo de dato & Valor Null & Primaria & For\'anea & Comentario \\ \hline
		UserIdUsuario & Cadena & Si & Si & No & \\ \hline
		RolFKIdRol & Entero & Si & No & Si & \\ \hline
		UserNombreUsuario & Cadena & Si & No & No & \\ \hline
		UserCorreo & Cadena & Si & No & No & \\ \hline
		UserNombres & Cadena & Si & No & No & \\ \hline
		UserApellidos & Cadena & Si & No & No & \\ \hline
	\end{tabular}
\end{table}

\begin{table}[ht]
	\caption{Tabla Semillero}
	\label{labelTableSemillero}
	\begin{tabular}{ |l|l|l|l|l|l| }
		\hline
		Nombre & Tipo de dato & Valor Null & Primaria & For\'anea & Comentario \\ \hline
		SmlrIdSemillero & Entero & Si & Si & No & \\ \hline 
		SmlrNombre & Cadena & Si & No & No & \\ \hline 
		SmlrDescripcion & Texto largo & Si & No & No & \\ \hline 
		SmlrRutaLogo & Texto largo & Si & No & No & \\ \hline 
		SmlrMision & Texto largo & Si & No & No & \\ \hline 
		SmlrVision & Texto largo & Si & No & No & \\ \hline 
		SmlrRuta & Cadena & Si & No & No & \\ \hline 
		SmlrSigla & Cadena & Si & No & No & \\ \hline 
		SmlrJustificacion & Texto largo & Si & No & No & \\ \hline 
		SmlrMetodologiaTrabajo & Texto largo & Si & No & No & \\ \hline 
		SmlrAreasTrabajo & Texto largo & Si & No & No & \\ \hline 	
	\end{tabular}
\end{table}

\begin{table}[ht]
	\caption{Tabla Integrantes Semillero}
	\label{labelTableIntegrantesSemillero}
	\begin{tabular}{ |l|l|l|l|l|l| }
		\hline
		Nombre & Tipo de dato & Valor Null & Primaria & For\'anea & Comentario \\ \hline
		IsmlIdIntegrante & Entero & Si & Si & No & \\ \hline 
		UserFKIdUsuario & Entero & Si & No & Si & \\ \hline 
		SmlrFKIdSemillero & Entero & Si & No & Si & \\ \hline 
		IsmlFecha & Fecha & Si & No & No & \\ \hline 
		IsmlEstado & Cadena & Si & No & No & \\ \hline	
	\end{tabular}
\end{table}

\begin{table}[ht]
	\caption{Tabla Rol}
	\label{labelTableRol}
	\begin{tabular}{ |l|l|l|l|l|l| }
		\hline
		Nombre & Tipo de dato & Valor Null & Primaria & For\'anea & Comentario \\ \hline
		RolIdRol & Entero & Si & Si & No & \\ \hline 
		RolNombreRol & Cadena & Si & No & No & \\ \hline	
	\end{tabular}
\end{table}

\begin{table}[ht]
	\caption{Tabla Permisos}
	\label{labelTablePermisos}
	\begin{tabular}{ |l|l|l|l|l|l| }
		\hline
		Nombre & Tipo de dato & Valor Null & Primaria & For\'anea & Comentario \\ \hline
		PrmsIdPermiso & Entero & Si & Si & No & \\ \hline 
		PrmsNombrePermiso & Cadena & Si & No & No & \\ \hline 
		PrmsIcono & Cadena & Si & No & No & \\ \hline 	
	\end{tabular}
\end{table}


\begin{table}[ht]
	\caption{Tabla Subpermisos}
	\label{labelTableSubpermisos}
	\begin{tabular}{ |l|l|l|l|l|l| }
		\hline
		Nombre & Tipo de dato & Valor Null & Primaria & For\'anea & Comentario \\ \hline
		SpmsIdSubpermiso & Entero & Si & Si & No & \\ \hline 
		SpmsNombreSubpermiso & Cadena & Si & No & No & \\ \hline 
		SpmsURL & Cadena & Si & No & No & \\ \hline 
		PrmsFKIdPermiso & Entero & Si & No & Si & \\ \hline 	
	\end{tabular}
\end{table}

\begin{table}[ht]
	\caption{Tabla Menú Usuario}
	\label{labelTableMenuUsuario}
	\begin{tabular}{ |l|l|l|l|l|l| }
		\hline
		Nombre & Tipo de dato & Valor Null & Primaria & For\'anea & Comentario \\ \hline
		RolFKIdRol & Entero & Si & Si & Si & \\ \hline 
		SpmsFKIdSubpermiso & Entero & Si & Si & Si & \\ \hline 
		MusrEstado & Cadena & Si & No & No & \\ \hline 	
	\end{tabular}
\end{table}

\begin{table}[ht]
	\caption{Tabla Repositorio}
	\label{labelTableRepositorio}
	\begin{tabular}{ |l|l|l|l|l|l| }
		\hline
		Nombre & Tipo de dato & Valor Null & Primaria & For\'anea & Comentario \\ \hline
		RepoIdRepositorio & Entero & Si & Si & No & \\ \hline 
		RepoNombre & Cadena & Si & No & No & \\ \hline 
		RepoDescripcion & Texto largo & Si & No & No & \\ \hline 
		RepoRuta & Texto largo & Si & No & No & \\ \hline 
		RepoEstado & Cadena & Si & No & No & \\ \hline 	
	\end{tabular}
\end{table}


\begin{table}[ht]
	\caption{Tabla Documento}
	\label{labelTableDocumento}
	\begin{tabular}{ |l|l|l|l|l|l| }
		\hline
		Nombre & Tipo de dato & Valor Null & Primaria & For\'anea & Comentario \\ \hline
		DocuIdDocumento & Entero & Si & Si & No & \\ \hline 
		DocuNombre & Cadena & Si & No & No & \\ \hline 
		DocuRuta & Texto largo & Si & No & No & \\ \hline 
		DocuFecha & Fecha & Si & No & No & \\ \hline 
		DocuEstado & Cadena & No & No & No & \\ \hline 
		UserFKIdUsuario & BI Entero & Si & No & Si & \\ \hline 
		RepoFKIdRepositorio & Entero & Si & No & Si & \\ \hline 	
	\end{tabular}
\end{table}

\begin{table}[ht]
	\caption{Tabla Correos Enviados}
	\label{labelTableCorreosEnviados}
	\begin{tabular}{ |l|l|l|l|l|l| }
		\hline
		Nombre & Tipo de dato & Valor Null & Primaria & For\'anea & Comentario \\ \hline
		CrnvIdCorreoEnviado & Entero & Si & Si & No & \\ \hline 
		CrnvTipo & Cadena & Si & No & No & \\ \hline 
		CrnvDescripcion & Texto largo & Si & No & No & \\ \hline 
		CrnvEstadoEnvio & VARCHAR(50) & Si & No & No & \\ \hline 
		CrnvFechaEnvio & Fecha & Si & No & No & \\ \hline 
		SmlrFKIdSemillero & Entero & Si & No & Si & \\ \hline 
		UserFKDestinatario & Cadena & Si & No & Si & \\ \hline 	
	\end{tabular}
\end{table}


\begin{table}[ht]
	\caption{Tabla Acceso Sistema}
	\label{labelTableAccesoSistema}
	\begin{tabular}{ |l|l|l|l|l|l| }
		\hline
		Nombre & Tipo de dato & Valor Null & Primaria & For\'anea & Comentario \\ \hline
		AstmIdAcceso & BI Entero & Si & Si & No & \\ \hline 
		AstmFecha & Fecha & Si & No & No & \\ \hline 
		AstmIP & Cadena & Si & No & No & \\ \hline 
		UserFKIdUsuario & BI Entero & Si & No & Si & \\ \hline 	
	\end{tabular}
\end{table}


\begin{table}[ht]
	\caption{Tabla Solicitudes}
	\label{labelTableSolicitudes}
	\begin{tabular}{ |l|l|l|l|l|l| }
		\hline
		Nombre & Tipo de dato & Valor Null & Primaria & For\'anea & Comentario \\ \hline
		SlctIdSolicitud & Entero & Si & Si & No & \\ \hline 
		CursIdCurso & Entero & Si & No & No & \\ \hline 
		UserFKEstudiante & BI Entero & Si & No & Si & \\ \hline 
		UserFKAdmin & BI Entero & Si & No & Si & \\ \hline 
		SlctFechaSolicitud & Fecha & Si & No & No & \\ \hline 
		SlctFechaRespuesta & Fecha & No & No & No & \\ \hline 
		SlctEstado & Cadena & No & No & No & \\ \hline 
		SlctComentario & Texto largo & Si & No & No & \\ \hline 	
	\end{tabular}
\end{table}


\begin{table}[ht]
	\caption{Tabla Grupo Investigaci\'on}
	\label{labelTableGrupoInvestigacion}
	\begin{tabular}{ |l|l|l|l|l|l| }
		\hline
		Nombre & Tipo de dato & Valor Null & Primaria & For\'anea & Comentario \\ \hline
		GrivIdGrupoInv & Entero & Si & Si & No & \\ \hline 
		GrivFechaCreacion & Fecha & Si & No & No & \\ \hline 
		GrivNombre & Cadena & Si & No & No & \\ \hline 
		GrivMision & Texto largo & Si & No & No & \\ \hline 
		GrivVision & Texto largo & Si & No & No & \\ \hline 	
	\end{tabular}
\end{table}


\begin{table}[ht]
	\caption{Tabla Programa}
	\label{labelTablePrograma}
	\begin{tabular}{ |l|l|l|l|l|l| }
		\hline
		Nombre & Tipo de dato & Valor Null & Primaria & For\'anea & Comentario \\ \hline
		ProgIdPrograma & Entero & Si & Si & No & \\ \hline 
		ProgNombre & Cadena & Si & No & No & \\ \hline 
		ProgCodigo & Cadena & No & No & No & \\ \hline 
		FacuFKIdFacultad & Entero & Si & No & Si & \\ \hline 	
	\end{tabular}
\end{table}


\begin{table}[ht]
	\caption{Tabla Programas Semillero}
	\label{labelTableProgramasSemillero}
	\begin{tabular}{ |l|l|l|l|l|l| }
		\hline
		Nombre & Tipo de dato & Valor Null & Primaria & For\'anea & Comentario \\ \hline
		ProgFKIdPrograma & Entero & Si & Si & Si & \\ \hline 
		SmlrFKIdSemillero & Entero & Si & Si & Si & \\ \hline 
		PgsmEstado & ENUM('T', 'F') & Si & No & No & \\ \hline 	
	\end{tabular}
\end{table}


\begin{table}[ht]
	\caption{Tabla Programas Grupos Inv}
	\label{labelTableProgramasGruposInv}
	\begin{tabular}{ |l|l|l|l|l|l| }
		\hline
		Nombre & Tipo de dato & Valor Null & Primaria & For\'anea & Comentario \\ \hline
		GrivFKIdGrupoInv & Entero & Si & Si & Si & \\ \hline 
		ProgFKIdPrograma & Entero & Si & Si & Si & \\ \hline 
		PggrEstado & ENUM('T', 'F') & Si & No & No & \\ \hline 
		PggrFecha & Fecha & Si & No & No & \\ \hline 
		PggrCartaAval & Texto largo & Si & No & No & \\ \hline 	
	\end{tabular}
\end{table}


\begin{table}[ht]
	\caption{Tabla Proyectos}
	\label{labelTableProyectos}
	\begin{tabular}{ |l|l|l|l|l|l| }
		\hline
		Nombre & Tipo de dato & Valor Null & Primaria & For\'anea & Comentario \\ \hline
		ProyIdProyecto & Entero & Si & Si & No & \\ \hline 
		ProyNombre & Texto largo & Si & No & No & \\ \hline 
		ProyDescripcion & Texto largo & Si & No & No & \\ \hline 
		ProyTipo & Cadena & Si & No & No & \\ \hline 
		ProSitado & Cadena & Si & No & No & \\ \hline 
		LinvFKIdLineaInv & Entero & Si & No & Si & \\ \hline 
		ConvFKIdConvocatoria & Entero & Si & No & Si & \\ \hline 
		GrivFKIdGrupoInv & Entero & Si & No & Si & \\ \hline 
		ProgFKIdPrograma & Entero & Si & No & Si & \\ \hline 
		ProyFormulacionProblema & Texto largo & Si & No & No & \\ \hline 
		ProyJustificacion & Texto largo & Si & No & No & \\ \hline 
		ProyMarcoTeorico & Texto largo & Si & No & No & \\ \hline 
		ProyMetodologia & Texto largo & Si & No & No & \\ \hline 
		ProyBibliografia & Texto largo & Si & No & No & \\ \hline 
		ProyAnexos & Texto largo & Si & No & No & \\ \hline 	
	\end{tabular}
\end{table}


\begin{table}[ht]
	\caption{Tabla Rol Semillero}
	\label{labelTableRolSemillero}
	\begin{tabular}{ |l|l|l|l|l|l| }
		\hline
		Nombre & Tipo de dato & Valor Null & Primaria & For\'anea & Comentario \\ \hline
		RsmlIdRolSemillero & Entero & Si & Si & No & \\ \hline 
		RsmlNombre & Cadena & Si & No & No & \\ \hline 
		RsmlEstado & Cadena & Si & No & No & \\ \hline 
		SmlrFKIdSemillero & Entero & Si & No & Si & \\ \hline 	
	\end{tabular}
\end{table}


\begin{table}[ht]
	\caption{Tabla Par Evaluador}
	\label{labelTableParEvaluador}
	\begin{tabular}{ |l|l|l|l|l|l| }
		\hline
		Nombre & Tipo de dato & Valor Null & Primaria & For\'anea & Comentario \\ \hline
		PaevIdParEvaluador & Entero & Si & Si & No & \\ \hline 
		PaevFecha & Fecha & Si & No & No & \\ \hline 
		PaevEstado & Cadena & Si & No & No & \\ \hline 
		PaevTipo & Cadena & Si & No & No & \\ \hline 
		UserFKParEvaluador & Cadena & Si & No & Si & \\ \hline 	
	\end{tabular}
\end{table}


\begin{table}[ht]
	\caption{Tabla Actividad Semillero}
	\label{labelTableActividadSemillero}
	\begin{tabular}{ |l|l|l|l|l|l| }
		\hline
		Nombre & Tipo de dato & Valor Null & Primaria & For\'anea & Comentario \\ \hline
		ActsIdActividad & Entero & Si & Si & No & \\ \hline 
		ActsDescripcion & Texto largo & Si & No & No & \\ \hline 
		ActsFechaCreacion & Fecha & Si & No & No & \\ \hline 
		ActsEstado & Cadena & Si & No & No & \\ \hline 
		ActsFechaInicio & DATE & Si & No & No & \\ \hline 
		ActsFechaFin & DATE & Si & No & No & \\ \hline 
		MetsFKIdMetaSemillero & Entero & Si & No & Si & \\ \hline 	
	\end{tabular}
\end{table}


\begin{table}[ht]
	\caption{Tabla Convocatoria}
	\label{labelTableConvocatoria}
	\begin{tabular}{ |l|l|l|l|l|l| }
		\hline
		Nombre & Tipo de dato & Valor Null & Primaria & For\'anea & Comentario \\ \hline
		ConvIdConvocatoria & Entero & Si & Si & No & \\ \hline 
		ConvNombre & Texto largo & Si & No & No & \\ \hline 
		ConvDescripcion & Texto largo & Si & No & No & \\ \hline 
		ConvObservaciones & Texto largo & No & No & No & \\ \hline 
		ConvFechaCreacion & Fecha & Si & No & No & \\ \hline 
		ConvFechaInicio & DATE & Si & No & No & \\ \hline 
		ConvFechaFin & DATE & Si & No & No & \\ \hline 
		ConvEstado & Cadena & Si & No & No & \\ \hline 
		TconFKIdTipoConvocatoria & Entero & Si & No & Si & \\ \hline 
		ConvOferente & Cadena & Si & No & No & \\ \hline 
		ConvDestinatario & Cadena & Si & No & No & \\ \hline 	
	\end{tabular}
\end{table}


\begin{table}[ht]
	\caption{Tabla Banco Preliminar}
	\label{labelTableBancopreliminar}
	\begin{tabular}{ |l|l|l|l|l|l| }
		\hline
		Nombre & Tipo de dato & Valor Null & Primaria & For\'anea & Comentario \\ \hline
		BpreIdBancoPreliminar & Entero & Si & Si & No & \\ \hline 
		BpreDescripcion & Cadena & Si & No & No & \\ \hline 
		BpreEstado & Cadena & Si & No & No & \\ \hline 
		ConvFKIdConvocatoria & Entero & Si & No & Si & \\ \hline 
		BpreFecha & Fecha & Si & No & No & \\ \hline 	
	\end{tabular}
\end{table}


\begin{table}[ht]
	\caption{Tabla Proyectos preliminares}
	\label{labelTableProyectospreliminares}
	\begin{tabular}{ |l|l|l|l|l|l| }
		\hline
		Nombre & Tipo de dato & Valor Null & Primaria & For\'anea & Comentario \\ \hline
		PpreIdProyectoPreliminar & Entero & Si & Si & No & \\ \hline 
		ProyFKIdProyecto & Entero & Si & No & Si & \\ \hline 
		BpreFKIdBancoPreliminar & Entero & Si & No & Si & \\ \hline 
		PpreEstado & Cadena & Si & No & No & \\ \hline 
		PpreFecha & Fecha & Si & No & No & \\ \hline 	
	\end{tabular}
\end{table}


\begin{table}[ht]
	\caption{Tabla Banco definitivo}
	\label{labelTableBancodefinitivo}
	\begin{tabular}{ |l|l|l|l|l|l| }
		\hline
		Nombre & Tipo de dato & Valor Null & Primaria & For\'anea & Comentario \\ \hline
		BdefIdBancoDefinitivo & Entero & Si & Si & No & \\ \hline 
		BdefEstado & Cadena & Si & No & No & \\ \hline 
		BdefFecha & Fecha & Si & No & No & \\ \hline 
		ConvFKIdConvocatoria & Entero & Si & No & Si & \\ \hline 	
	\end{tabular}
\end{table}


\begin{table}[ht]
	\caption{Tabla Proyectos definitivos}
	\label{labelTableProyectosdefinitivos}
	\begin{tabular}{ |l|l|l|l|l|l| }
		\hline
		Nombre & Tipo de dato & Valor Null & Primaria & For\'anea & Comentario \\ \hline
		PdefIdProyectoDefinitivo & Entero & Si & Si & No & \\ \hline 
		PdefEstado & Cadena & Si & No & No & \\ \hline 
		PdefFecha & Fecha & Si & No & No & \\ \hline 
		BdefFKIdBancoDefinitivo & Entero & Si & No & Si & \\ \hline 
		ProyFKIdProyecto & Entero & Si & No & Si & \\ \hline 	
	\end{tabular}
\end{table}


\begin{table}[ht]
	\caption{Tabla Evaluaci\'on}
	\label{labelTableEvaluacion}
	\begin{tabular}{ |l|l|l|l|l|l| }
		\hline
		Nombre & Tipo de dato & Valor Null & Primaria & For\'anea & Comentario \\ \hline
		EvlcIdEvaluacion & BI Entero & Si & Si & No & \\ \hline 
		EvlcDescripcion & Cadena & Si & No & No & \\ \hline 
		EvlcFecha & Fecha & Si & No & No & \\ \hline 
		EvlcEstado & Cadena & Si & No & No & \\ \hline 
		EvlcResultado & Cadena & Si & No & No & \\ \hline 
		AparFKIdAsignacion & Entero & Si & No & Si & \\ \hline 	
	\end{tabular}
\end{table}


\begin{table}[ht]
	\caption{Tabla Tipo Convocatoria}
	\label{labelTableTipoConvocatoria}
	\begin{tabular}{ |l|l|l|l|l|l| }
		\hline
		Nombre & Tipo de dato & Valor Null & Primaria & For\'anea & Comentario \\ \hline
		TconIdTipoConvocatoria & Entero & Si & Si & No & \\ \hline 
		TconNombre & Cadena & Si & No & No & \\ \hline 
		TconEstado & Cadena & Si & No & No & \\ \hline 	
	\end{tabular}
\end{table}


\begin{table}[ht]
	\caption{Tabla Linea Investigaci\'on}
	\label{labelTableLineaInvestigacion}
	\begin{tabular}{ |l|l|l|l|l|l| }
		\hline
		Nombre & Tipo de dato & Valor Null & Primaria & For\'anea & Comentario \\ \hline
		LinvIdLineaInv & Entero & Si & Si & No & \\ \hline 
		LinvNombreLinea & Cadena & Si & No & No & \\ \hline 
		LinvDescripcion & Texto largo & Si & No & No & \\ \hline 
		ProgFKIdPrograma & Entero & Si & No & Si & \\ \hline 	
	\end{tabular}
\end{table}


\begin{table}[ht]
	\caption{Tabla Facultad}
	\label{labelTableFacultad}
	\begin{tabular}{ |l|l|l|l|l|l| }
		\hline
		Nombre & Tipo de dato & Valor Null & Primaria & For\'anea & Comentario \\ \hline
		FacuIdFacultad & Entero & Si & Si & No & \\ \hline 
		FacuNombre & Cadena & Si & No & No & \\ \hline 
		FacuEstado & Cadena & Si & No & No & \\ \hline 	
	\end{tabular}
\end{table}


\begin{table}[ht]
	\caption{Tabla Evaluador Interno}
	\label{labelTableEvaluadorInterno}
	\begin{tabular}{ |l|l|l|l|l|l| }
		\hline
		Nombre & Tipo de dato & Valor Null & Primaria & For\'anea & Comentario \\ \hline
		EvaiIdEvaluadorInterno & Entero & Si & Si & No & \\ \hline 
		ProgFKIdPrograma & Entero & Si & No & Si & \\ \hline 
		PaevFKIdParEvaluador & Entero & Si & No & Si & \\ \hline 	
	\end{tabular}
\end{table}


\begin{table}[ht]
	\caption{Tabla Asignaci\'on Par Evaluaci\'on}
	\label{labelTableAsignacionParEvaluacion}
	\begin{tabular}{ |l|l|l|l|l|l| }
		\hline
		Nombre & Tipo de dato & Valor Null & Primaria & For\'anea & Comentario \\ \hline
		AparIdAsignacion & Entero & Si & Si & No & \\ \hline 
		ProyFKIdProyecto & Entero & Si & No & Si & \\ \hline 
		PaevFKIdParEvaluador & Entero & Si & No & Si & \\ \hline 
		AparEstado & Cadena & Si & No & No & \\ \hline 
		AparFecha & Fecha & Si & No & No & \\ \hline 
		UserFKAsignadoPor & Cadena & Si & No & Si & \\ \hline 	
	\end{tabular}
\end{table}


\begin{table}[ht]
	\caption{Tabla Preevaluacion}
	\label{labelTablePreevaluacion}
	\begin{tabular}{ |l|l|l|l|l|l| }
		\hline
		Nombre & Tipo de dato & Valor Null & Primaria & For\'anea & Comentario \\ \hline
		PrevIdPreevaluacion & Entero & Si & Si & No & \\ \hline 
		PrevDescripcion & Texto largo & Si & No & No & \\ \hline 
		PrevFecha & Fecha & Si & No & No & \\ \hline 
		PrevEstado & Cadena & Si & No & No & \\ \hline 
		PrevResultado & Cadena & Si & No & No & \\ \hline 
		InstFKIdInstancia & Entero & Si & No & Si & \\ \hline 
		ProyFKIdProyecto & Entero & Si & No & Si & \\ \hline 
		UserFKIdUsuario & Cadena & Si & No & Si & \\ \hline 	
	\end{tabular}
\end{table}


\begin{table}[ht]
	\caption{Tabla Instancia}
	\label{labelTableInstancia}
	\begin{tabular}{ |l|l|l|l|l|l| }
		\hline
		Nombre & Tipo de dato & Valor Null & Primaria & For\'anea & Comentario \\ \hline
		InstIdInstancia & Entero & Si & Si & No & \\ \hline 
		InstNombre & Cadena & Si & No & No & \\ \hline 
		InstDescripcion & Texto largo & Si & No & No & \\ \hline 
		InstEstado & Cadena & Si & No & No & \\ \hline 	
	\end{tabular}
\end{table}


\begin{table}[ht]
	\caption{Tabla Integrantes Grupo}
	\label{labelTableIntegrantesGrupo}
	\begin{tabular}{ |l|l|l|l|l|l| }
		\hline
		Nombre & Tipo de dato & Valor Null & Primaria & For\'anea & Comentario \\ \hline
		IgruIdIntegrante & Entero & Si & Si & No & \\ \hline 
		UserFKIdUsuario & Cadena & Si & No & Si & \\ \hline 
		GrivFKIdGrupoInv & Entero & Si & No & Si & \\ \hline 
		IgruEstado & Cadena & Si & No & No & \\ \hline 
		IgruFecha & Fecha & Si & No & No & \\ \hline 	
	\end{tabular}
\end{table}


\begin{table}[ht]
	\caption{Tabla Proyectos Semilleros}
	\label{labelTableProyectosSemilleros}
	\begin{tabular}{ |l|l|l|l|l|l| }
		\hline
		Nombre & Tipo de dato & Valor Null & Primaria & For\'anea & Comentario \\ \hline
		PrsmIdProyecto & Entero & Si & Si & No & \\ \hline 
		PrsmFecha & Fecha & Si & No & No & \\ \hline 
		PrsmEstado & Cadena & Si & No & No & \\ \hline 
		ProyFKIdProyecto & Entero & Si & No & Si & \\ \hline 
		SmlrFKIdSemillero & Entero & Si & No & Si & \\ \hline 	
	\end{tabular}
\end{table}


\begin{table}[ht]
	\caption{Tabla Proyectos Grupos}
	\label{labelTableProyectosGrupos}
	\begin{tabular}{ |l|l|l|l|l|l| }
		\hline
		Nombre & Tipo de dato & Valor Null & Primaria & For\'anea & Comentario \\ \hline
		PrgrIdProyecto & Entero & Si & Si & No & \\ \hline 
		PrgrFecha & Fecha & Si & No & No & \\ \hline 
		PrgrEstado & Cadena & Si & No & No & \\ \hline 
		ProyFKIdProyecto & Entero & Si & No & Si & \\ \hline 
		GrivFKIdGrupoInv & Entero & Si & No & Si & \\ \hline 
		PrgrLugarEjecucion & Cadena & Si & No & No & \\ \hline 
		MpioFKIdMpio & Entero & Si & No & Si & \\ \hline 
		PrgrDuracion & DOUBLE & Si & No & No & \\ \hline 
		TinvFKIdTipo & Entero & Si & No & Si & \\ \hline 
		PrgrValorTotal & BI Entero & Si & No & No & \\ \hline 
		PrgrValorSolicitado & BI Entero & Si & No & No & \\ \hline 
		PrgrValorContrapartida & BI Entero & Si & No & No & \\ \hline 
		PrgrTotalInvestigadores & Entero & Si & No & No & \\ \hline 
		PrgrTipo & Cadena & Si & No & No & \\ \hline 
		PrgrResultadosEsperados & Texto largo & Si & No & No & \\ \hline 
		PrgrEstraComunicacion & Texto largo & Si & No & No & \\ \hline 
		PrgrPalabrasClave & Texto largo & Si & No & No & \\ \hline 
		& Cadena & No & No & No & \\ \hline 	
	\end{tabular}
\end{table}


\begin{table}[ht]
	\caption{Tabla Roles Integrantes Semillero}
	\label{labelTableRolesIntegrantesSemillero}
	\begin{tabular}{ |l|l|l|l|l|l| }
		\hline
		Nombre & Tipo de dato & Valor Null & Primaria & For\'anea & Comentario \\ \hline
		RismIdRol & Entero & Si & Si & No & \\ \hline 
		RismFecha & Fecha & Si & No & No & \\ \hline 
		RismEstado & Cadena & Si & No & No & \\ \hline 
		IsmlFKIdIntegrante & Entero & Si & No & Si & \\ \hline 
		RsmlFKIdRolSemillero & Entero & Si & No & Si & \\ \hline 	
	\end{tabular}
\end{table}


\begin{table}[ht]
	\caption{Tabla Rol Grupo}
	\label{labelTableRolGrupo}
	\begin{tabular}{ |l|l|l|l|l|l| }
		\hline
		Nombre & Tipo de dato & Valor Null & Primaria & For\'anea & Comentario \\ \hline
		RgruIdRolGrupo & Entero & Si & Si & No & \\ \hline 
		RgruNombre & Cadena & Si & No & No & \\ \hline 
		RgruEstado & Cadena & Si & No & No & \\ \hline 
		GrivFKIdGrupoInv & Entero & Si & No & Si & \\ \hline 	
	\end{tabular}
\end{table}


\begin{table}[ht]
	\caption{Tabla Roles Integrantes Grupo}
	\label{labelTableRolesIntegrantesGrupo}
	\begin{tabular}{ |l|l|l|l|l|l| }
		\hline
		Nombre & Tipo de dato & Valor Null & Primaria & For\'anea & Comentario \\ \hline
		RigrIdRol & Entero & Si & Si & No & \\ \hline 
		RigrFecha & Fecha & Si & No & No & \\ \hline 
		RigrEstado & Cadena & Si & No & No & \\ \hline 
		IgruFKIdIntegrante & Entero & Si & No & Si  & \\ \hline 
		RgruFKIdRolGrupo & Entero & Si & No & Si & \\ \hline 	
	\end{tabular}
\end{table}


\begin{table}[ht]
	\caption{Tabla Objetivos Proyecto Semillero}
	\label{labelTableObjetivosProyectoSemillero}
	\begin{tabular}{ |l|l|l|l|l|l| }
		\hline
		Nombre & Tipo de dato & Valor Null & Primaria & For\'anea & Comentario \\ \hline
		ObpsIdObjetivo & Entero & Si & Si & No & \\ \hline 
		ObpsID & Cadena & Si & No & No & \\ \hline 
		ObpsDescripcion & Cadena & Si & No & No & \\ \hline 
		ObpsFecha & Fecha & Si & No & No & \\ \hline 
		ObpsEstado & Cadena & Si & No & No & \\ \hline 
		ObpsTipo & Cadena & Si & No & No & \\ \hline 
		PrsmFKIdProyecto & Entero & Si & No & Si & \\ \hline 	
	\end{tabular}
\end{table}


\begin{table}[ht]
	\caption{Tabla Informe Parcial}
	\label{labelTableInformeParcial}
	\begin{tabular}{ |l|l|l|l|l|l| }
		\hline
		Nombre & Tipo de dato & Valor Null & Primaria & For\'anea & Comentario \\ \hline
		InpaIdInformeParcial & Entero & Si & Si & No & \\ \hline 
		InpaDescripcion & Texto largo & Si & No & No & \\ \hline 
		InpaFecha & Fecha & Si & No & No & \\ \hline 
		InpaEstado & Cadena & Si & No & No & \\ \hline 
		InpaMetodologia & Texto largo & Si & No & No & \\ \hline 
		InpaResultados & Texto largo & Si & No & No & \\ \hline 
		InpaImpacto & Texto largo & Si & No & No & \\ \hline 
		InpaEstraDivulgacion & Texto largo & Si & No & No & \\ \hline 
		ProyFKIdProyecto & Entero & Si & No & Si & \\ \hline 
		InpaDocSoporte & Texto largo & Si & No & No & \\ \hline 
		InpaPresupuesto & Texto largo & Si & No & No & \\ \hline 	
	\end{tabular}
\end{table}


\begin{table}[ht]
	\caption{Tabla Integrantes Proyecto Semillero}
	\label{labelTableIntegrantesProyectoSemillero}
	\begin{tabular}{ |l|l|l|l|l|l| }
		\hline
		Nombre & Tipo de dato & Valor Null & Primaria & For\'anea & Comentario \\ \hline
		IprsIdIntegrante & Entero & Si & Si & No & \\ \hline 
		IprsFecha & Fecha & Si & No & No & \\ \hline 
		IprsEstado & Cadena & Si & No & No & \\ \hline 
		RolpFKIdRolProyecto & Entero & Si & No & Si & \\ \hline 
		IsmlFKIdIntegrante & Entero & Si & No & Si & \\ \hline 
		PrsmFKIdProyecto & Entero & Si & No & Si &  \\ \hline 	
	\end{tabular}
\end{table}


\begin{table}[ht]
	\caption{Tabla Rol Proyecto}
	\label{labelTableRolProyecto}
	\begin{tabular}{ |l|l|l|l|l|l| }
		\hline
		Nombre & Tipo de dato & Valor Null & Primaria & For\'anea & Comentario \\ \hline
		RolpIdRolProyecto & Entero & Si & Si & No & \\ \hline 
		RolpNombre & Cadena & Si & No & No & \\ \hline 
		RolpEstado & Cadena & Si & No & No & \\ \hline 	
	\end{tabular}
\end{table}


\begin{table}[ht]
	\caption{Tabla Informe Final}
	\label{labelTableInformeFinal}
	\begin{tabular}{ |l|l|l|l|l|l| }
		\hline
		Nombre & Tipo de dato & Valor Null & Primaria & For\'anea & Comentario \\ \hline
		InfiIdInformeFinal & Entero & Si & Si & No & \\ \hline 
		InfiProblemaPlanteado & Texto largo & Si & No & No & \\ \hline 
		InfiHipotesis & Texto largo & Si & No & No & \\ \hline 
		InfiAvanceObtenido & Texto largo & Si & No & No & \\ \hline 
		InfiRetosPlanteados & Texto largo & Si & No & No & \\ \hline 
		InfiGradoComprobacion & Texto largo & Si & No & No & \\ \hline 
		InfiGestionProyecto & Texto largo & Si & No & No & \\ \hline 
		InfiObservaciones & Texto largo & Si & No & No & \\ \hline 
		InfiRecomendaciones & Texto largo & Si & No & No & \\ \hline 
		InfiActores & Texto largo & Si & No & No & \\ \hline 	
	\end{tabular}
\end{table}


\begin{table}[ht]
	\caption{Tabla Informe Final}
	\label{labelTableInformeFinal}
	\begin{tabular}{ |l|l|l|l|l|l| }
		\hline
		Nombre & Tipo de dato & Valor Null & Primaria & For\'anea & Comentario \\ \hline
		InfiIdInformeFinal & Entero & Si & Si & No & \\ \hline 
		InfiProblemaPlanteado & Texto largo & Si & No & No & \\ \hline 
		InfiHipotesis & Texto largo & Si & No & No & \\ \hline 
		InfiAvanceObtenido & Texto largo & Si & No & No & \\ \hline 
		InfiRetosPlanteados & Texto largo & Si & No & No & \\ \hline 
		InfiGradoComprobacion & Texto largo & Si & No & No & \\ \hline 
		InfiGestionProyecto & Texto largo & Si & No & No & \\ \hline 
		InfiObservaciones & Texto largo & Si & No & No & \\ \hline 
		InfiRecomendaciones & Texto largo & Si & No & No & \\ \hline 
		InfiActores & Texto largo & Si & No & No & \\ \hline 	
	\end{tabular}
\end{table}


\begin{table}[ht]
	\caption{Tabla Cronograma Convocatoria}
	\label{labelTableCronogramaConvocatoria}
	\begin{tabular}{ |l|l|l|l|l|l| }
		\hline
		Nombre & Tipo de dato & Valor Null & Primaria & For\'anea & Comentario \\ \hline
		CrcvIdCronograma & Entero & Si & Si & No & \\ \hline 
		CrcvDescripcion & Texto largo & Si & No & No & \\ \hline 
		CrcvFechaCreacion & Fecha & Si & No & No & \\ \hline 
		CrcvFechaActividad & Fecha & Si & No & No & \\ \hline 
		CrcvEstado & Cadena & No & No & No & \\ \hline 
		ConvFKIdConvocatoria & Entero & Si & No & Si & \\ \hline 	
	\end{tabular}
\end{table}

\begin{table}[ht]
	\caption{Tabla Soportes Cronograma}
	\label{labelTableSoportesCronograma}
	\begin{tabular}{ |l|l|l|l|l|l| }
		\hline
		Nombre & Tipo de dato & Valor Null & Primaria & For\'anea & Comentario \\ \hline
		SocrIdSoporte & Entero & Si & Si & No & \\ \hline 
		SocrDescripcion & Texto largo & Si & No & No & \\ \hline 
		SocrRutaAdjunto & Cadena & Si & No & No & \\ \hline 
		SocrEstado & Cadena & Si & No & No & \\ \hline 
		SocrFecha & Fecha & Si & No & No & \\ \hline 
		CrcvFKIdCronograma & Entero & Si & No & Si &  \\ \hline	
	\end{tabular}
\end{table}

\begin{table}[ht]
	\caption{Tabla Producci\'on Grupo}
	\label{labelTableProduccionGrupo}
	\begin{tabular}{ |l|l|l|l|l|l| }
		\hline
		Nombre & Tipo de dato & Valor Null & Primaria & For\'anea & Comentario \\ \hline
		PgruIdProduccion & Entero & Si & Si & No & \\ \hline 
		Pgru & Cadena & No & No & No & \\ \hline 	
	\end{tabular}
\end{table}

\begin{table}[ht]
	\caption{Tabla Unidades Acad\'emicas}
	\label{labelTableUnidadesAcademicas}
	\begin{tabular}{ |l|l|l|l|l|l| }
		\hline
		Nombre & Tipo de dato & Valor Null & Primaria & For\'anea & Comentario \\ \hline		
		UnacIdUnidad & Entero & Si & Si & No & \\ \hline 
		UnacNombre & Cadena & Si & No & No & \\ \hline 
		UnacEstado & Cadena & Si & No & No & \\ \hline 
		UnacFechaCreacion & Fecha & Si & No & No & \\ \hline	
	\end{tabular}
\end{table}
	
	% Finalizaci�n
	%
%\makeglossaries
%\loadglsentries{glossaires}
%
%\newglossaryentry{latex}
%{
%	name=latex,
%	description={Is a mark up language specially suited 
%		for scientific documents}
%}
%
%\newglossaryentry{maths}
%{
%	name=mathematics,
%	description={Mathematics is what mathematicians do}
%}

%\title{How to create a glossary}
%\author{ }
%\date{ }

%\maketitle

%The \Gls{latex} typesetting markup language is specially suitable 
%for documents that include \gls{maths}. 

%\clearpage

%\printglossaries
	\section[CONCLUSIONES]{CONCLUSIONES}
	% \section[Bibliograf�a]{Bibliograf�a}
	
	\bibliographystyle{alpha}
	\bibliography{referencias}
	
	\section[ANEXOS]{ANEXOS}
\end{document}